\chapter[General Conclusion and Perspectives]{General Conclusion and\\ Perspectives}
\label{ch:master-conclusion}

\section{General Conclusion}

Recovering a Software Bill of Materials from a compiled Android application is
not only a detection problem. It is a problem of deciding what can be claimed
from incomplete, transformed, and sometimes conflicting evidence. A component
name or exact version may look precise while depending on weak namespace
overlap, an incomplete reference collection, a feature-equivalent adjacent
release, or source information that does not prove binary presence. Trust
therefore depends on how evidence is qualified and bounded, not merely on
whether a tool returns an answer.

This thesis studied that problem as a research question in Android
third-party-library detection and binary \SBOM reconstruction. The first
chapter introduced the context, motivation, research questions, and scope of
the work. Chapter~\ref{ch:master-background} then established the technical
background needed to understand software supply chains, \SBOM standards,
vulnerability identifiers, Android compilation, APK evidence, and the
methodological constraints of a structured literature review. It showed why a
compiled APK cannot be treated as a transparent dependency manifest and why
binary evidence must be interpreted with provenance and uncertainty.

Chapter~\ref{ch:master-sota} organized the state of the art by recovery role,
evidence source, transformation exposure, version capability, validation
practice, and trust implication. The reviewed literature provides many useful
signals: package and class profiles, structural graphs, fuzzy signatures,
opcode and literal features, API references, learned representations,
optimization-aware matching, and benchmark-driven assessment. However, the
synthesis also showed recurring limitations. Exact-version identification is
fragile under shrinking and optimization; reported scores are often tied to
selected datasets; source-known dependencies are not always equivalent to
binary presence; and most approaches stop before producing standard,
provenance-aware \SBOM output.

The central conclusion is therefore conditional. Android APK evidence can
support useful reconstructed component claims, but only when the claim remains
as precise as the evidence allows. Library identity and exact version must be
distinguished. Missing reference coverage, near-tie versions, transformed
code, native payloads, and dynamic loading must be visible rather than hidden.
Vulnerability interpretation must inherit the confidence of the component and
version evidence instead of bypassing it. In this view, abstention and bounded
claims are not failures; they are necessary protections against false
precision.

\section{Answers to the Research Questions}

Table~\ref{tab:master-rq-answers} summarizes the answers to the supporting
research questions introduced in Chapter~\ref{ch:master-introduction}. The
table is retained because it gives a compact link between the literature
review and the final synthesis.

\begingroup
\small
\setlength{\LTpre}{0.5em}
\setlength{\LTpost}{0.5em}
\begin{longtable}{>{\raggedright\arraybackslash}p{0.07\textwidth}>{\raggedright\arraybackslash}p{0.29\textwidth}>{\raggedright\arraybackslash}p{0.52\textwidth}}
\caption{Answers to the supporting research questions.}
\label{tab:master-rq-answers}\\
\toprule
\textbf{RQ} & \textbf{Question focus} & \textbf{Answer}\\
\midrule
\endfirsthead
\toprule
\textbf{RQ} & \textbf{Question focus} & \textbf{Answer}\\
\midrule
\endhead
RQ1 & Evidence sources available in compiled Android applications &
Recoverable evidence includes archive metadata, manifest information, DEX
packages and classes, methods, opcodes, literals, API references, dependency
and call relations, resources, native-file observations, reference artifacts,
and external validation sources. These sources differ in directness, version
specificity, resilience, and authority.\\

RQ2 & How existing approaches combine evidence & Existing approaches range
from namespace and class-set matching to structural graphs, opcode and
semantic features, fuzzy signatures, learned representations, metadata
coherence, and multi-stage reference matching. The literature supports
multi-source reasoning, but no combination is universally reliable under
shrinking, optimization, repackaging, feature-equivalent releases, and
incomplete references.\\

RQ3 & Factors that determine identity and version strength & Evidence strength
depends on directness, uniqueness, independence, transformation resilience,
version specificity, provenance reliability, reference coverage, and
reproducibility. Identity and exact-version support must be assessed
separately because a signal can identify a library family while remaining
unable to distinguish adjacent releases.\\

RQ4 & Representation of conflict, absence, transformation, and source
dependence & These conditions should be represented through competing
candidates, bounded version sets, conflict records, coverage limits,
transformation exposure, report-only outcomes, and abstention. Missing
evidence should not be interpreted as component absence, and source-known
dependencies should not automatically be treated as proof of APK-side
presence.\\

RQ5 & Information required for an auditable and security-relevant claim & A
claim should identify the subject APK, candidate component and artifact family,
claim type and granularity, version or version set, supporting and opposing
evidence, provenance, coverage, transformation exposure, confidence class,
ambiguity, validation context, and permitted downstream use. Vulnerability
interpretation must inherit these boundaries.\\

RQ6 & Evaluation without overstating accuracy & Evaluation should separate
controlled fixtures, clean source-validated cases, partial source-validated
cases, source-known upper bounds, diagnostics, report-only experiments, and
unsupported limitations. It should publish denominators, distinguish emitted
row correctness from full dependency recall, retain negative experiments and
abstentions, and avoid promoting signals without control evidence.\\
\bottomrule
\end{longtable}
\endgroup

Together, these answers shift evaluation from the question ``Did the detector
emit a version?'' to the stronger question ``What evidence supports this
claim, what could invalidate it, and how may it safely be used?''

\section{Future Works and Perspectives}

\begingroup
\small
\setlength{\LTpre}{0.5em}
\setlength{\LTpost}{0.5em}
\begin{longtable}{>{\raggedright\arraybackslash}p{0.25\textwidth}>{\raggedright\arraybackslash}p{0.35\textwidth}>{\raggedright\arraybackslash}p{0.28\textwidth}}
\caption{Priority perspectives for future Android APK-to-SBOM research.}
\label{tab:master-future-directions}\\
\toprule
\textbf{Direction} & \textbf{Research objective} & \textbf{Required evidence}\\
\midrule
\endfirsthead
\toprule
\textbf{Direction} & \textbf{Research objective} & \textbf{Required evidence}\\
\midrule
\endhead
Empirical calibration & Calibrate identity, exact-version, ambiguity, and
abstention classes by component family and transformation profile & Held-out
ground truth, reliability analysis, and error-sensitive calibration metrics.\\
Independent validation & Apply recovery and reporting rules to other tools,
analysts, and datasets & External tools, blinded case review, and inter-rater
agreement on claim classes.\\
Public benchmark design & Measure exactness, coverage, abstention,
provenance, and vulnerability-authority correctness together & Reproducible
APKs, complete denominators, transformation labels, and
version-discriminating ground truth.\\
Provenance-aware \SBOM profiles & Standardize how reconstructed claims expose
evidence, ambiguity, and validation context & Interoperable CycloneDX or
related profiles tested across producers and consumers.\\
Native and dynamic attribution & Extend analysis to shared libraries,
protected payloads, and downloaded code & Separate corpora,
architecture-aware fingerprints, runtime provenance, and scope-specific
metrics.\\
Human factors and explainability & Determine whether analysts understand and
act appropriately on claim classes, conflicts, and abstentions & User studies,
decision-quality measures, and explanation comparisons.\\
\bottomrule
\end{longtable}
\endgroup

Future work should first improve benchmark design. A useful benchmark should
not only ask whether a detector found a name. It should record exact
versions, transformed variants, reference coverage, native and dynamic
boundaries, and the denominator used for each metric. Such a benchmark would
make it possible to compare tools without confusing emitted-row precision
with full dependency recall.

A second perspective is calibrated uncertainty. Instead of returning one
forced exact version, future systems should support bounded version sets,
confidence classes, and abstention policies whose empirical behavior is
measured on independent data. This is especially important for adjacent
versions, optimized code, and incomplete reference catalogs.

A third direction is interoperability. Reconstructed \SBOM producers need a
shared vocabulary for observed binary evidence, inferred identity, candidate
version sets, external validation, and unsupported regions. CycloneDX
extensions or profiles could preserve these distinctions while remaining
usable by vulnerability-management tools.

Finally, Android research should expand beyond Maven bytecode. Native
libraries, dynamically loaded code, protected payloads, and generated
artifacts require their own evidence models and validation protocols. Treating
them as ordinary Maven components would create false certainty. Treating them
as explicit future research domains keeps the boundary honest while opening
the path toward broader mobile software transparency.
